%
%      Brno University of Technology
%      Faculty of Information Technology
%
%      MSc Thesis	2009/2010
%
%      Design of S-Boxes Using Genetic Algorithms
%
%
%      Bedrich Hovorka
%
%
%
Evolutionary algorithms are informed methods for searching state spaces,
which \cite{evo} are based on the metaphor of evolution in nature. The solution of a problem is
converted into a process of evolving a population of randomly generated solutions. Each
solution is encoded into a string of symbols (parameters) and evaluated by a so-called fitness
function, which expresses the quality of the solution -- the higher the value, the more
promising the given solution is and the more frequently it enters the reproduction process, during which
new solutions are generated. The population of solutions is commonly called a population of individuals
or chromosomes. The reproduction process is based on two driving forces:

\begin{itemize}
\item variation operators crossover and mutation, which bring diversity
to the population
\item selection, which favors high-quality individuals
\end{itemize}

The combination of variation and selection generally contributes to improving the fitness function of individuals
in the newly forming population. In the process of crossing individuals, just as in living nature,
new individuals/offspring are obtained by crossing mostly pairs of parental
individuals. All components of the evolutionary process are stochastic -- for example, pairs with better fitness function
enter the reproduction process more frequently, but weaker
individuals also have a chance to enter the reproduction process.

First, the basic concepts are shown using the most widespread type of evolutionary algorithm. For other
algorithms, only the differences are listed.

\section{Genetic Algorithm}\label{secGA}

The chromosome (individual) encoding the solution is represented by a binary vector
(string) of constant length $n$:
\begin{description}
  \item $\vec{X} = (X_0, X_1, \dots, X_{n-1})$, where $X_i$ is the $i$-th variable of the string
  \item $\vec{x} = (x_0, x_1, \dots, x_{n-1})$ is the string of specific instances of variables
  \item $X_i = x_i, x_i \in \{0,1\}$
  \item $D = (X^1, X^2,\dots, X^N), X^j \in D$ is a set of $N$ strings that specifies population $D$
  \item $D \subseteq \{0,1\}^n$.
\end{description}

Let $f$ be an objective/cost function defined over the set of binary vectors
of length $n$
\begin{equation}
  f: \{0,1\}^n \rightarrow \mathcal{R}
\end{equation}
which evaluates each binary vector $\vec{x}$ with a real number. The goal is to find the global
extremum of function $f$. In the case of a minimization problem, this means finding the vector
\begin{equation}
\vec{x}_{opt} = \argmin\limits_{\vec{x} \in \{0,1\}^n} f(\vec{x})
\end{equation}
Function $f$ is usually transformed into a fitness function $F$ such that
the original optimization problem is converted into a maximization problem and
an appropriate scale of the fitness function is achieved. Using this transformation
also facilitates changing the selection pressure, which significantly affects the convergence
of the evolutionary process. The operation of the standard GA algorithm is described by algorithm \ref{geneticAlgorithmCode}.

\begin{algoritmus}{Pseudocode of genetic algorithm}{geneticAlgorithmCode}
  \STATE Set $t=0$, randomly generate initial population $D(0)$ with cardinality $N$,
  \WHILE{termination condition for the algorithm is not satisfied}
    \STATE Perform evaluation of individuals in population $D(t)$ with fitness function $F(X)$,
    \STATE Generate offspring population $O(t)$ with cardinality $M \le N$ using crossover
and mutation operators,
    \STATE Create new population $D(t+1)$ by replacing part of population $D(t)$ with individuals from
$O(t)$,
    \STATE Set $t \leftarrow t+1$,
  \ENDWHILE
\end{algoritmus}

In the reproduction process, pairs of parent individuals are usually selected quasi-randomly
for subsequent crossover and mutation. Part of the individuals from the offspring
population is then used to replace individuals in the original population $D(t)$, and thus
we obtain a new population $D(t+1)$. Termination of the evolutionary process is either given
by the maximum number of generations or by detection of stagnation of the evolutionary process.
A genetic algorithm can have several types of crossover and mutation. Some operators are designed
according to the type of problem being solved. Mutation is usually performed by inverting bit(s) at random positions.
Crossover can be, for example, one-point or two-point. After selecting two parents, first a
point(s) is chosen at which the chromosome of both parents is divided into two parts. Two offspring are created.
The first of them is composed of the first part of the first parent and the second part of the second parent.
And the chromosome of the second offspring is assembled from the remaining parts. Selection can be performed by:
\begin{itemize}
  \item tournament among $N$ randomly selected individuals,
  \item proportionally (the ratio of individual fitness to the sum of fitness of all individuals) also known as roulette wheel
  \item linear ranking (selection probability is a linear function of position after sorting population by fitness)
  \item exponential ranking (like linear, but probability depends exponentially)
\end{itemize}


\section{Genetic Programming}
\begin{description}
\item[Representation] Genetic programming \cite{ui4} works with so-called executable structures.
Most often, these
are programs that are represented by trees and graphs. These structures have
(but do not always have to have) variable length. Unlike classical GA, here we work directly
with an executable structure that represents both the genotype and phenotype of the candidate solution.

\item[Genetic operators] work on executable structures. In addition to common
operators (such as crossover and mutation), there is a number of advanced operators
enabling the generation of programs with subroutines, modules, etc.
\item[Evaluation] As part of determining the fitness value, the code of the candidate program
(generally an executable structure) is executed for a defined set of inputs and
the obtained results are evaluated.
\end{description}


While in general genetic programming there is a search in the solution space that
changes during evolution, in this text we will only deal with static state space.
In genetic programming, one can design any executable structure that best
captures the problem being solved. And design custom operators for this structure.


\subsection{Cartesian Genetic Programming}\label{algorithmCGP}
Cartesian Genetic Programming (CGP) \cite{vasicek} can be understood as a variant of genetic programming
that uses an acyclic graph with a fixed number of nodes in the representation of the problem being solved.
Each of the nodes can implement one of several fixed functions and can be connected to the output of one of the previous nodes.
Connectivity is influenced by the so-called \emph{l-back} parameter.
The chromosome is a certain sequence of integer values having a fixed length. In this, one can
see the division into genotype and phenotype as in classical GA.
The algorithm used in CGP corresponds to the evolutionary strategy $(1+L)$,
where the value of $L$ is on the order of units.
Only the mutation operator is used.

\subsubsection{Encoding the problem}
The chromosome is composed of $m \cdot n$ triplets ($m$ - number of columns, $n$ - number of rows) of integer values, where the first two values indicate the output number to which the first and second input of the element corresponding to the triplet is connected, the third value indicates which logical function the element implements.
At the end of the chromosome is an $o$-tuple of size equal to the number of outputs of the circuit being searched.
Each element of the $o$-tuple determines the index of the element (its output)
to which the primary output of the circuit will be connected. Other parameters are $i$ - number of inputs,
$l$ - \emph{l-back}, $F$ - set of functions.

Consider the following form of chromosome,
corresponding to a simple Boolean function\footnote{By eliminating gates from it, one can obtain AND,
simply the simplest bent function given for illustration},
satisfying SAC:
\begin{verbatim}
    {2,1,2,2,2,1,3}([2]1,1,8)([3]0,0,7)([4]1,3,0)([5]3,2,8)(5)
\end{verbatim}
The form of the chromosome is explained in appendix \ref{cgpfunctions}.
In figure \ref{simpleBent}, this chromosome is depicted in graph form.

\insertimage[260px]{simpleBent}{Example of a simple Boolean function satisfying SAC found by
Cartesian genetic programming}
\subsubsection{Program implementation}
During \cite{millerThomson} generation of new individuals or during mutation, several conditions must be satisfied so that
the chromosome represents a correct program.
Let ${c^n}_{kj}$ represent the gene of the $k$-th input in
column $j$ (the leftmost column has $j = 0$), then it must satisfy these constraints:
\begin{eqnarray}
  c_{kj} < e_{\max} & if & j < l \\
  e_{\min} \leq c_{kj} < e_{\max} & if & j \geq l \\
  e_{\min} = i + (j-l) n \\
  e_{\max} = i + j n
\end{eqnarray}
Let ${c^o}_k$ represent the gene of the $k$-th output of the entire circuit, then for it
these constraints apply:
\begin{eqnarray}
  h_{\min} \leq {c^o}_k < h_{\max} \\
  h_{\min} = i + (m - l) n \\
  h_{\max} = i + (m - 1) n
\end{eqnarray}
Let ${c^f}_k$ represent the function gene of the $k$-th node, then:
\begin{equation}
  {c^f}_k \in F
\end{equation}
From these conditions, it is also evident that the implementation of crossover,
in which it would be necessary to comply with these conditions,
would be more complicated, and therefore it is preferably not used.


\section{EDA}\label{algorithmEDA}
\cite{evo} During the crossover operation of a classical genetic algorithm, building blocks are broken up.
This fact is a frequent cause of poor convergence of the optimization process.
Therefore, algorithms have been designed that try to solve this problem by taking a global view
of the entire population. In the entire population, statistics of occurrences and dependencies of alleles at
individual positions are computed. According to these statistics, new individuals are subsequently generated.
These algorithms are collectively called EDA -- Estimation of Distribution Algorithm.

\subsection{Operation of EDA algorithm}
Unlike GA, in each step $M < N$ \uv{promising} individuals are selected, from
which an estimate of the probability distribution of selected individuals $p(X | Ds(t))$ is performed.
Then a set of offspring $O(t)$ with cardinality $M < N$ is generated (by sampling the obtained distribution).
A new population $D(t+1)$ is created by replacing part of population $D(t)$ with individuals from $O(t)$.


The fundamental problem is the estimation of probability distribution
$p(X|Ds(t))$. From theory, the relationship for joint probability distribution
of a random vector using local probability distributions
with the use of conditional probabilities is known:
\begin{equation}
p(X) = \prod\limits^{n-1}_{i=0} p(X_i | X_0 , X_1,\dots, X_{i-1})
\end{equation}
And it is precisely in this computation (complexity of the probability model) that EDA algorithms are divided
into several types, which are listed below.

\subsection{UMDA}\label{algorithmUMDA}
UMDA (Univariate Marginal Distribution Algorithm) uses the simplest
statistical model. It assumes that loci are mutually independent, the occurrence of alleles
at each position is governed by a distribution function.
Let population $P$ of size $N$ consist of chromosomes of length $n$. For each
locus $i \in \{0 .. n-1\}$ and each allele $x_i \in \{0,\dots,r_{i-1}\}$ we define $p_i(x_i)$ as the frequency
of occurrence of strings that have $x_i$ at the $i$-th position in set $P$.
\begin{equation}
  p_i(x_i) = \frac{n_i(x_i)}{N}
\end{equation}
The probability with which a new individual
$X = a_1 a_2 \dots a_{n-1}$ will be generated by the UMDA algorithm must then be
\begin{equation}
  p(X) = \prod\limits^{n-1}_{i=0} p_i(a_i)
\end{equation}
Thus, first the frequencies of different alleles are computed at each position
from the population of promising solutions, then for each new individual the bit at the $i$-th position
is set to value $a$ with probability $p_i(a)$. The dependency graph corresponding to
the UMDA algorithm does not contain any edges ($E = \emptyset$).

\subsection{BMDA}\label{algorithmBMDA}
In addition to independent loci, the BMDA (Bivariate Marginal Distribution
Algorithm) algorithm also allows pairwise dependencies between loci, which are detected
by statistical tests. For generating values at dependent positions, the apparatus
of conditional probability is used.
Again, we consider population $P$ of size $N$ with chromosomes of length $n$. For any two
positions $i \neq j \in \{0..n-1\}$ and any two possible alleles $x_i \in \{0,\dots,r_{i-1}\}$,
$x_j \in \{0,\dots,r_{j-1}\}$ at
these positions, we define $p_{i,j}(x_i,x_j)$ as the frequency of occurrence of strings with values $x_i$
and $x_j$ at positions $i$ and $j$:
\begin{equation}
  p_{i,j} = \frac{n_{i,j}(i,j)}{N}
\end{equation}
The conditional probability
\begin{equation}
  p_{i,j}(x_i | x_j) = \frac{p_{i,j}(x_i, x_j)}{p_j(x_j)}
\end{equation}
is then estimated from the frequency of occurrence of allele $x_i$ at position $i$ in cases when $x_j$ is
at the $j$-th position.
The dependency graph \cite{evo, pelikan} for the BMDA algorithm has a tree structure, each gene is
dependent on at most one other gene. Some groups of genes may be
completely independent, then the graph has the form of a set of trees.
The probability of generating an individual $X = a_1 a_2 \dots a_{n-1}$ is then
\begin{equation}
p(X) = \prod\limits_{j \in R} p(a_j) \prod\limits_{(i,j) \in E} p(a_j|a_i)
\end{equation}
where $R$ resp. $E$ is the set of root nodes (i.e., independent variables) resp. edges of graph $G$.

For testing the independence of qualitative attributes in statistics, the $\chi^2$ distribution is used
\begin{equation}
  \chi^2 = \sum\frac{real\_value - expected\_value}{expected\_value}
\end{equation}
In our case, this involves testing the independence of contents of two loci:
\begin{equation}
  \chi^2_{i,j} = \sum\limits_{x_i = 0}^{r_i -1} \sum\limits_{x_j = 0}^{r_j -1}
  \frac{(Np_{i,j}(x_i, x_j) - Np_i(x_i)p_j(x_j))^2}{Np_i(x_i)p_j(x_j)}
\end{equation}
where $r_i$ is the number of different alleles at locus $i$, $r_j$ is the number of different alleles at locus $j$ (for
binary chromosomes $r = s = 2$). This formula is based on the assumption that for
the occurrence of values at independent positions, according to probability theory,
$p_{i, j} (x_i , x_j) = p_i(x_i) p_j(x_j)$ must hold.
For all pairs of loci, contingency tables are constructed during population traversal.
Using these tables, a set of dependent pairs is then created using the $\chi^2$-test, which
is ordered according to the value of $\chi^2$.
Two binary loci are dependent at 95\% confidence if the $\chi^2$ value is greater than 3.84.
From the set of dependent pairs, a dependency tree is formed, which is then used to generate
new individuals.

%algoritmy vytvareni graf a generovani novych jedincu jen pokud bude malo textu
Other variants of EDA algorithms (e.g., BOA) were not used for the reasons stated below.

\section{Multi-criteria optimization}
The task of searching for S-boxes can be based on multiple criteria that may conflict with each other.
In multi-criteria optimization, we do not search for a simple optimum as in single-criterion optimization, but we try to find a suitable
compromise when optimizing individual criteria. Formally, we search for a parameter vector
$ \vec{x}^* = [x^*_1, x^*_2, \dots, x^*_n]^T$ that:
\begin{itemize}
  \item satisfies $m$ constraints: $g_i(\vec{x}) \leq 0, i = 1,2,\dots, m$
  \item satisfies $p$ constraints: $h_i(\vec{x}) = 0, i = 1,2,\dots, p$
  \item will optimize the function: $f(\vec{x}) = \left(f_1(\vec{x}), f_2(\vec{x}), \dots, f_p(\vec{x}) \right)$
\end{itemize}
For this reason, it was necessary to incorporate this problem into evolutionary algorithms.
Several approaches have been designed for this purpose.

\subsection{Pareto optimum}
\begin{itemize}
  \item For any two vectors $\vec{a}$, $\vec{b}$, we say that vector $\vec{a}$ strongly
  dominates vector $\vec{b}$ if:
  \begin{enumerate}
    \item $f_i(\vec{a}) \geq f_i(\vec{b})$ for $\forall i = 1,\dots, p$
    \item and simultaneously: $f_i(\vec{a}) > f_i(\vec{b})$ for at least one $i$
  \end{enumerate}
  \item Vector $\vec{a}$ then weakly dominates $\vec{b}$ if only the first condition holds
  \item Based on Pareto dominance, the concept of \uv{Pareto Optimum} can be defined:
  Vector $\vec{a}$ is Pareto optimal if there does not exist any vector $\vec{b}$ such that it dominates vector $\vec{a}$.
  In other words: vector $\vec{a}$ is Pareto optimal if in the decision
  space there does not exist any vector $\vec{b}$ that would reduce some criterion without simultaneously increasing at least one other criterion.
  \item However, this concept does not give us just a single solution, but usually an entire set of solutions,
  the so-called Pareto optimal solution set (Pareto optimalset).
  \item Vectors $\vec{a}$ that are also in the Pareto optimal set (POS) are
called non-dominated. The curve/hypersurface passing through these non-dominated vectors from POS is
called the Pareto front.
  \item The solution to MOP is then any point from the Pareto front
\end{itemize}

\subsection{VEGA algorithm}\label{algorithmVEGA}
VEGA is the simplest
possible multi-criteria GA. For several criteria (denoted $M$), the
division of the GA population in each generation into $M$ equal subgroups must be managed randomly.
Each subgroup is assigned a fitness function based on a criterion function,
different for each subgroup. The population is divided into $M$ equal parts in each generation. Each
individual in the first subgroup is assigned a fitness function based only on the first
criterion function, while each individual in the second subgroup is assigned
a fitness function according to the second criterion function, and so on. It is better to shuffle
the population before it is divided into subgroups. For each solution,
the assigned fitness function, selection operator, is reserved among solutions of each
subgroup. The selection operator is applied until the complete
subgroup is filled. From that point on, all individuals in the subgroup have assigned
fitness functions based on a single criterion function. This restricts the selection
operation only within the subgroup and emphasizes good solutions corresponding to that
particular criterion function. In VEGA, a proportional selection
operator is used in the subpopulations.

\subsection{SPEA algorithm}\label{algorithmSPEA}
This algorithm \cite{evo, thiele} maintains an external elite population $P_e$. This population contains a fixed
number of non-dominated solutions that have been found since the beginning of evolution.
In each generation, newly found non-dominated solutions are compared with
the existing external population and the resulting non-dominating solutions are retained.
SPEA does more than just preserve the elite. It also uses this elite to participate in
genetic operations together with the current population.

This algorithm starts with a randomly created population $P_0$ of size $N$ and
an empty external population ${P_0}_e$ with maximum capacity $N_e$. In each generation $t$,
the best non-dominated solutions (belonging to the first non-dominated front) from population
$P_t$ are copied to the external population ${P_t}_e$. Dominated solutions in the modified external
population are found and removed from this population.
In the external population remain the best non-dominated solutions of the combined
population containing old and new elites. In order to limit population overgrowth,
the size of the external population is bound to limit $N_e$. If the population size
equals the external population size, all solutions are also contained in the external
population. However, if the population size exceeds $N_e$, not all elites
can be placed in the external population.
Elite individuals who are in less crowded regions in the non-dominated front
are stored. Once the new elites are preserved for the next generation, the algorithm
turns to the current population and using genetic operators, finds a new
population.
The first step is to assign fitness to each solution in the population (including external).
In fact, in the SPEA algorithm, fitness (strength) $S_i$ is first assigned to
each member $i$ from the external population. The strength $S_i$ is proportional to the number of current
population members ($n_i$) who dominate the external solution $i$:
\begin{equation}
  S_i = \frac{n_i}{N+1}
\end{equation}
This ensures that greater strength is given to the elite that dominates more solutions in the current
population. Division by the value $(N + 1)$ ensures that the maximum strength value
of any member of the external population will always be less than $1$. Moreover, a non-dominated
solution dominating smaller solutions has better fitness. The fitness
of the current population member $j$ is determined as the sum of strengths of all external members that
strongly\footnote{in the original text \cite{evo} it is weakly, but in the implementation strong dominance was used}
dominate solution $j$:
\begin{equation}
  F_j = 1 + \sum\limits_{i \in {P_t}_e \wedge i \leq j} S_i
\end{equation}
Adding one makes the fitness value of any current population member $P_t$
greater than the fitness value of any external population member ${P_t}_e$.
External population members receive smaller fitness values than members of the new
population. Population
members who are dominated by many external members receive a large
fitness value. Therefore, during selection over both populations, smaller
values are preferred (so-called minimization). Both populations are then used as a whole during crossover
and mutation.

\subsubsection{Clustering}\label{secClustering}
The clustering algorithm reduces the size of the external population ${P_t}_e$ from ${N'}_e$
to $N_e$ (where ${N'}_e > N_e$). Initially, each solution in ${P_t}_e$ is considered as one solution
in one separate cluster. Thus, there are initially ${N'}_e$ clusters in the external population.
Subsequently, distances between individual pairs of clusters are computed. Typically,
the distance $d_{12}$ between two clusters $C_1$ and $C_2$ is defined as the average
Euclidean distance of all pairs of solutions:
\begin{equation}
 d_{12} = \dfrac{\sum\limits_{i \in C_1 , j \in C_2} d(i,j) }{|C_1| |C_2|}
\end{equation}
Once all clusters are computed, the two clusters with minimum
distance are merged together and jointly form one larger cluster. This
merging process continues until the number of clusters in the external population is reduced
to size $N_e$. Then in each cluster, only the solution with minimum
average distance from other solutions in the cluster is kept, and all other solutions are
deleted.

\begin{table}[bt]
   \begin{center}
    \begin{longtable}{| r || p{0.4\textwidth} | p{0.4\textwidth} | }
      \hline
      ~ & Advantages & Disadvantages \\ \hline\hline
      VEGA
      & \begin{itemize}
          \item simple idea, easy to implement
          \item requires only minor changes to single-criterion GA
          \item same computational complexity as GA
          \item tendency to find solutions close to the individual best solution of each criterion
        \end{itemize}
      & \begin{itemize}
          \item evaluates each solution with only a single criterion function
          \item does not find diverse solutions in the population, directs all solutions toward one champion
        \end{itemize}
      \\ \hline
      SPEA
      & \begin{itemize}
         \item stores solutions closest to the Pareto-optimal front throughout the entire evolution
         \item clustering allows better distribution of the limited size of optimal solutions
        \end{itemize}
      & \begin{itemize}
          \item parameter $N_e$ -- balance with $N$ must be ensured by an appropriate ratio
        \end{itemize}
      \\ \hline
    \end{longtable}
    \caption{Comparison of algorithms for multi-criteria optimization}
    \label{vegaSpeaComparison}
  \end{center}
\end{table}

In table \ref{vegaSpeaComparison}, both algorithms are summarized. While VEGA favors one
criterion more, SPEA tries to find a broader spectrum of solutions.

\pagebreak

\section{Other methods for searching S-boxes}
Other state space search methods could also be used for searching S-boxes.
For informed methods, the criterion is used as a heuristic function of the given method.
The goal state is defined by the required value of the criterion, time, or
number of steps (generations). The use of heuristics in S-box search is mentioned, for example,
in \cite{goodsboxes}.

\subsubsection{Gradient algorithm -- hill climbing}
Selects \cite{ui1} for expansion the derived node that is most promising. Both parent and siblings are forgotten.
If all derived nodes are worse, the algorithm terminates. Thus it often ends in a local extremum.
\subsubsection{Ordered search}
Unlike the gradient algorithm, it stores unexpanded nodes in a priority queue
(or sorted list) according to evaluation. The disadvantage is high memory requirements.
The queue size can be limited, but this will make the search approach gradient search and
it will be prone to termination/getting stuck in a local extremum.
\subsubsection{Random search}
Used in \cite{twofish}. This search is the only blind one among those mentioned.
It derives a new solution by a random change (mutation) of the current solution.
The best solution can be stored continuously. It also serves for comparison when tuning the evolutionary algorithm.
\subsubsection{Parallel random search}
Unlike random search, which is in principle completely blind,
this search generates in one step multiple mutants from the current individual
(works with a population of solutions -- parallelism) and selects from them the best one
that differs from the current one, and it is used for mutations in the next step.
De facto, it is used in Cartesian Genetic Programming
(see page \pageref{algorithmCGP}).

A comparison of random and parallel random search
is part of the following text in subsection \ref{secParalelRandomComparison}.

\section{S-box search vs. symbolic regression}
If S-box search, or at least part of it, can be
converted to some theoretical problem solvable by means of
evolution, it will most likely be symbolic regression.
Symbolic regression means finding a mathematical expression that describes data from a training
set. The expression can be represented as an individual in the evolutionary process, for which
during the computation of the evaluation function, values from the training set are tested.
The evaluation function is given, for example, by the sum of squares of deviations at required outputs or
the number of matches.
The evolution process is in this view analogous to neural network learning.
The expression can be represented as a tree (original genetic programming),
a grammar derivation (grammatical evolution), or any instance of an executable structure
or chromosome that best captures the problem being solved,
if, for example, we are looking for an optimal hardware
or software implementation.

When searching for S-boxes, we can consider these search variants:
\begin{enumerate}
  \item There is no training/testing set to speak of.
    The search is performed according to security criteria, therefore it is not symbolic regression.
  \item An expression is sought that exactly satisfies the training set, obtained from a previous phase.
  For example, if we first find a suitable (secure) permutation
  and then let an optimal implementation be found for it.
  The training set consists of all possible input combinations.
\end{enumerate}

Symbolic regression can only be used as an auxiliary mechanism for finding the optimal
implementation. However, symbolic regression must first find an exact implementation.
